\documentclass[12pt]{jlreq}
\usepackage{amsmath, amssymb}

\newcommand{\C}{\mathbb{C}}
\newcommand{\R}{\mathbb{R}}
\newcommand{\Z}{\mathbb{Z}}
\newcommand{\N}{\mathbb{N}}
\newcommand{\F}{\mathbb{F}}
\newcommand{\Prob}{\mathbb{P}}
\newcommand{\Exp}{\mathbb{E}}
\newcommand{\dd}{\,d}
\newcommand{\Ker}{\operatorname{Ker}}
\newcommand{\Impart}{\operatorname{Im}}
\newcommand{\Hom}{\operatorname{Hom}}

\begin{document}

\begin{enumerate}
  \item 次の条件をすべてみたす区間 \([0,1]\) 上の実数値ルベーグ可測関数 \(f(x)\) の例を構成せよ。
  \begin{enumerate}
    \item \(\displaystyle\int_0^1 |f(x)|\dd x<\infty\)。
    \item 任意の正の実数 \(a\le1\) に対し，
    \[
      \int_0^a |f(x)|^2\dd x=\infty.
    \]
  \end{enumerate}
  \item 次の条件をすべてみたす実数軸上の実数値ルベーグ可測関数 \(g(x)\) の例を構成せよ。
  \begin{enumerate}
    \item \(\displaystyle\int_{-\infty}^{\infty}|g(x)|\dd x<\infty\)。
    \item 任意の実数 \(a,b\) で \(a<b\) となるものに対し，
    \[
      \int_a^b |g(x)|^2\dd x=\infty.
    \]
    \item 集合 \(\{x\mid g(x)\ne0\}\) のルベーグ測度は \(1\) 以下である。
  \end{enumerate}
\end{enumerate}

\end{document}